% =============================================================================
% Capítulo 5: Scripts del Datawarehouse
% Entregable 2 del Avance N°1: script SQL de creación del modelo multidimensional.
% Las cifras de la Tabla de verificación se obtuvieron del catálogo de
% PostgreSQL el 05-11-2026 (ver docs/bitacora.md).
% =============================================================================
\section{Scripts del Datawarehouse}
\label{sec:scripts}

Este capítulo presenta los scripts que crean y verifican el Datawarehouse. En
el Avance N°1 el script crea la estructura del modelo multidimensional
\emph{sin datos}; los scripts de carga se agregan en el Avance N°2.

\subsection{Scripts Entregados}

\begin{table}[H]
  \centering
  \caption{Scripts del Datawarehouse y de apoyo}
  \label{tab:scripts}
  \small
  \begin{tabularx}{\linewidth}{@{}>{\raggedright\arraybackslash}p{5.4cm}
                                  >{\raggedright\arraybackslash}X@{}}
    \toprule
    \textbf{Archivo} & \textbf{Propósito} \\
    \midrule
    \bd{sql/01_ddl.sql} & Crea el esquema \bd{dwh} con las seis dimensiones,
      la tabla de hechos, sus restricciones, índices y comentarios \\
    \bd{sql/99_validacion.sql} & Consulta el catálogo de PostgreSQL para
      verificar tablas, columnas, llaves, restricciones, índices y
      nomenclatura \\
    \bd{tests/test_ddl.py} & 29 pruebas automáticas (pytest) que crean una
      base de prueba, ejecutan el script dos veces y comparan el resultado con
      el diseño \\
    \bd{docker-compose.yml} & Define los servicios PostgreSQL 17 y pgAdmin 4 \\
    \bd{herramientas/capturar_pgadmin.py} & Ejecuta el script y la validación
      en pgAdmin y toma las capturas del informe (Playwright) \\
    \bd{herramientas/generar_diccionario.py} & Genera el diccionario de datos
      del Anexo~A desde el catálogo de la base \\
    \bottomrule
  \end{tabularx}
\end{table}

\subsection{Script de Creación del Modelo}

El script \bd{sql/01_ddl.sql} está escrito para PostgreSQL 17 y se ejecuta de
corrido, sin intervención, tanto con \bd{psql} como en el Query Tool de
pgAdmin. Su estructura es la siguiente:

\begin{enumerate}[leftmargin=*]
  \item \textbf{Codificación:} \bd{SET client_encoding = 'UTF8'}. El archivo
        está en UTF-8, y esta línea asegura que las tildes de los comentarios
        y de las reglas se guarden bien aunque el script se ejecute con
        \bd{psql} desde una consola de Windows, que usa otra codificación.
  \item \textbf{Inicio de la transacción:} \bd{BEGIN}, seguido de
        \bd{DROP SCHEMA IF EXISTS dwh CASCADE} y \bd{CREATE SCHEMA dwh}. El
        script es \emph{idempotente}: se puede ejecutar las veces que sea
        necesario y siempre deja el mismo resultado.
  \item \textbf{Dimensiones:} las seis tablas \bd{dim_}, cada una con su llave
        primaria, su llave natural única y sus reglas \bd{CHECK}.
  \item \textbf{Tabla de hechos:} \bd{fact_venta}, con sus siete llaves
        foráneas (reglas \bd{ON DELETE RESTRICT} y \bd{ON UPDATE RESTRICT}) y
        las reglas de sus medidas.
  \item \textbf{Comentarios:} \bd{COMMENT ON} en el esquema, las 7 tablas y
        las 58 columnas (66 comentarios), que documentan el modelo dentro de la
        propia base.
  \item \textbf{Índices:} uno por cada llave foránea de \bd{fact_venta}, porque
        PostgreSQL no los crea por su cuenta y aceleran los cruces con las
        dimensiones.
  \item \textbf{Fin de la transacción:} \bd{COMMIT}. Si cualquier sentencia
        falla, la transacción se revierte completa y no queda nada a medias.
\end{enumerate}

La Figura~\ref{fig:ddl_fact} muestra la creación de la tabla de hechos. El
script completo se entrega en el repositorio del proyecto.

\begin{figure}[H]
  \centering
  \caption{Fragmento de \bd{sql/01_ddl.sql}: creación de la tabla de hechos \bd{fact_venta}}
  \label{fig:ddl_fact}
  \sinunicode{\lstinputlisting[style=sql, firstline=212, lastline=255,
                   firstnumber=212]{../sql/01_ddl.sql}}
\end{figure}

La Figura~\ref{fig:ddl_ejecutado} muestra la ejecución del script en el Query
Tool de pgAdmin, y la Figura~\ref{fig:estructura_fact} la estructura de la
tabla de hechos resultante.

\captura{02_pgadmin_ddl_ejecutado.png}{Ejecución de \bd{sql/01_ddl.sql} en el Query Tool de pgAdmin\label{fig:ddl_ejecutado}}
\captura{03_estructura_fact_venta.png}{Estructura de la tabla de hechos \bd{fact_venta} en pgAdmin\label{fig:estructura_fact}}

\subsection{Reglas de Integridad}

Además de las llaves, el script declara 20 restricciones \bd{CHECK} que la
base de datos hace cumplir en cada carga. Así, el ETL no puede dejar datos
inconsistentes en el Datawarehouse:

\begin{itemize}[leftmargin=*]
  \item \textbf{Montos:} el monto neto es igual al bruto menos el descuento, y
        el descuento está entre cero y el monto bruto.
  \item \textbf{Medidas no negativas:} la cantidad es de al menos una unidad;
        distancia, duración y peso no pueden ser negativos, y la edad del
        cliente está entre 0 y 120 años.
  \item \textbf{Dimensión tiempo:} la llave AAAAMMDD y cada atributo derivado
        (día, día de la semana, fin de semana, semana del año, mes, trimestre,
        semestre y año) deben coincidir con la fecha.
  \item \textbf{Dominios cerrados:} el tipo de cliente (persona o empresa), el
        tipo de servicio (pasajeros o carga) y la línea de negocio, que solo
        acepta los seis nombres del enunciado.
\end{itemize}

\subsection{Entorno de Ejecución}

El Datawarehouse corre en PostgreSQL 17.11 dentro de un contenedor Docker,
junto con pgAdmin 4 (Figura~\ref{fig:docker}). El esquema estrella creado por
el script se visualiza en el ERD Tool de pgAdmin
(Figura~\ref{fig:erd}).

\captura{01_docker_servicios_activos.png}{Salida de \bd{docker compose ps}: PostgreSQL 17 y pgAdmin activos\label{fig:docker}}
\captura{04_erd_esquema_estrella.png}{Esquema estrella en el ERD Tool de pgAdmin\label{fig:erd}}

\subsection{Verificación del Datawarehouse}
\label{sec:verificacion}

El resultado del script se verificó de tres formas: con las consultas de
\bd{sql/99_validacion.sql}, con las 29 pruebas automáticas de
\bd{tests/test_ddl.py} y ejecutando el script dos veces seguidas, para
comprobar que es idempotente. La Tabla~\ref{tab:verificacion} resume los
elementos del esquema \bd{dwh} consultados en el catálogo de PostgreSQL.

\begin{table}[H]
  \centering
  \caption{Elementos del esquema \bd{dwh} verificados en el catálogo}
  \label{tab:verificacion}
  \small
  \begin{tabularx}{\linewidth}{@{}>{\raggedright\arraybackslash}X cc@{}}
    \toprule
    \textbf{Elemento} & \textbf{Esperado según el diseño} & \textbf{Obtenido} \\
    \midrule
    Tablas (1 de hechos y 6 dimensiones)    & 7  & 7  \\
    Columnas                                & 58 & 58 \\
    Llaves primarias                        & 7  & 7  \\
    Restricciones \bd{UNIQUE}               & 7  & 7  \\
    Llaves foráneas (todas en \bd{fact_venta}) & 7  & 7  \\
    Restricciones \bd{CHECK}                & 20 & 20 \\
    Índices sobre llaves foráneas           & 7  & 7  \\
    Comentarios (esquema, tablas y columnas) & 66 & 66 \\
    Restricciones o índices fuera de la norma de nomenclatura & 0 & 0 \\
    Filas cargadas (Avance N°1, sin datos)  & 0  & 0  \\
    \bottomrule
  \end{tabularx}
  \nota{Consultas ejecutadas en PostgreSQL 17.11 el 05-11-2026.}
\end{table}

Las pruebas automáticas crean una base de prueba aparte, ejecutan el script
dos veces y comparan la base resultante con el diseño de
\bd{docs/modelo_dimensional.md}: tablas, columnas, tipos, llaves, reglas e
índices. También intentan insertar datos inválidos, como un monto neto que no
cuadra con el bruto o una venta con un cliente inexistente, y comprueban que la
base los rechaza. Las 29 pruebas se aprobaron.

Las figuras siguientes son resultados de \bd{sql/99_validacion.sql}. La
Figura~\ref{fig:val_tablas} muestra las siete tablas del esquema con su
comentario. La Figura~\ref{fig:val_fk} muestra las siete llaves foráneas de
\bd{fact_venta}, todas con nombre normado y regla \bd{RESTRICT}. La
Figura~\ref{fig:val_resumen} confirma la estrella: solo la tabla de hechos
tiene llaves foráneas salientes, y \bd{dim_zona} es referenciada dos veces
(origen y destino). La Figura~\ref{fig:val_checks} lista las reglas de
integridad, y la Figura~\ref{fig:val_nomenclatura} muestra que ninguna
restricción ni índice queda fuera de la norma de nomenclatura.

\captura{05a_validacion_tablas.png}{\bd{sql/99_validacion.sql}: tablas del esquema \bd{dwh}\label{fig:val_tablas}}
\captura{05b_validacion_fk.png}{\bd{sql/99_validacion.sql}: llaves foráneas de la tabla de hechos\label{fig:val_fk}}
\captura{05c_validacion_resumen.png}{\bd{sql/99_validacion.sql}: resumen de columnas y llaves por tabla\label{fig:val_resumen}}
\captura{05d_validacion_checks.png}{\bd{sql/99_validacion.sql}: restricciones \bd{CHECK}\label{fig:val_checks}}
\captura{05e_validacion_nomenclatura.png}{\bd{sql/99_validacion.sql}: control de nomenclatura (0 filas fuera de norma)\label{fig:val_nomenclatura}}
